\documentclass[12pt]{article}
\usepackage[utf8]{inputenc}
\usepackage[T1]{fontenc}
\usepackage[spanish,es-nodecimaldot]{babel}
\usepackage{amsmath,amssymb,mathtools,bm}
\usepackage[a4paper,margin=1.75cm,top=1.55cm,bottom=1.65cm]{geometry}
\usepackage{microtype}
\setlength{\parindent}{0pt}
\setlength{\parskip}{0.42em}
\setlength{\abovedisplayskip}{6pt}
\setlength{\belowdisplayskip}{6pt}
\newcommand{\encabezado}[3]{%
{\Large\bfseries #1\par}\vspace{2pt}
{\normalsize #2\hfill #3\par}\vspace{6pt}\hrule\vspace{8pt}}
\begin{document}

\encabezado{Tarea 2 --- Mecánica de Fluidos}{Gustavo Gatica, José Rosas, Nicolás Ulloa}{Lunes 27 de abril de 2026}

\textbf{Pregunta 1 (3 pts).} Demuestre que el gradiente vectorial $\nabla\bm u$ de un campo
\[
\bm u=u_r\hat r+u_\theta\hat\theta+u_\phi\hat\phi
\]
en coordenadas esféricas es
\begin{align}
\nabla\bm u={}&
\frac{\partial u_r}{\partial r}\hat r\hat r
+\frac1r\left(\frac{\partial u_r}{\partial\theta}-u_\theta\right)\hat\theta\hat r
+\frac{1}{r\sin\theta}\left(\frac{\partial u_r}{\partial\phi}-u_\phi\sin\theta\right)\hat\phi\hat r\\
&+\frac{\partial u_\theta}{\partial r}\hat r\hat\theta
+\frac1r\left(\frac{\partial u_\theta}{\partial\theta}+u_r\right)\hat\theta\hat\theta
+\frac{1}{r\sin\theta}\left(\frac{\partial u_\theta}{\partial\phi}-u_\phi\cos\theta\right)\hat\phi\hat\theta\\
&+\frac{\partial u_\phi}{\partial r}\hat r\hat\phi
+\frac1r\frac{\partial u_\phi}{\partial\theta}\hat\theta\hat\phi
+\frac{1}{r\sin\theta}\left(\frac{\partial u_\phi}{\partial\phi}+u_r\sin\theta+u_\theta\cos\theta\right)\hat\phi\hat\phi.
\end{align}
Usando este resultado, escriba explícitamente las componentes de la ecuación de Euler en coordenadas esféricas.

\textit{Solución.} Introducimos las derivadas direccionales asociadas a la base ortonormal esférica,
\begin{equation}
D_r=\frac{\partial}{\partial r},\qquad
D_\theta=\frac1r\frac{\partial}{\partial\theta},\qquad
D_\phi=\frac{1}{r\sin\theta}\frac{\partial}{\partial\phi}.
\end{equation}
Con la convención diádica usada en el enunciado,
\begin{equation}
\nabla\bm u=\hat r\,D_r\bm u+\hat\theta\,D_\theta\bm u+\hat\phi\,D_\phi\bm u.
\end{equation}
Las derivadas de la base que se necesitan son
\begin{align}
&\partial_r\hat r=\partial_r\hat\theta=\partial_r\hat\phi=0,\\
&\partial_\theta\hat r=\hat\theta,\qquad
\partial_\theta\hat\theta=-\hat r,\qquad
\partial_\theta\hat\phi=0,\\
&\partial_\phi\hat r=\sin\theta\,\hat\phi,\qquad
\partial_\phi\hat\theta=\cos\theta\,\hat\phi,\qquad
\partial_\phi\hat\phi=-\sin\theta\,\hat r-\cos\theta\,\hat\theta.
\end{align}
Aplicando la regla del producto,
\begin{align}
D_r\bm u={}&
\frac{\partial u_r}{\partial r}\hat r
+\frac{\partial u_\theta}{\partial r}\hat\theta
+\frac{\partial u_\phi}{\partial r}\hat\phi,\\
D_\theta\bm u={}&\frac1r\left[
\left(\frac{\partial u_r}{\partial\theta}-u_\theta\right)\hat r
+\left(\frac{\partial u_\theta}{\partial\theta}+u_r\right)\hat\theta
+\frac{\partial u_\phi}{\partial\theta}\hat\phi\right],\\
D_\phi\bm u={}&\frac{1}{r\sin\theta}\left[
\left(\frac{\partial u_r}{\partial\phi}-u_\phi\sin\theta\right)\hat r
+\left(\frac{\partial u_\theta}{\partial\phi}-u_\phi\cos\theta\right)\hat\theta\right.\\
&\left.\hspace{3.3cm}+\left(\frac{\partial u_\phi}{\partial\phi}+u_r\sin\theta+u_\theta\cos\theta\right)\hat\phi\right].
\end{align}
Al sustituir estas tres expresiones en $\nabla\bm u=\hat r D_r\bm u+\hat\theta D_\theta\bm u+\hat\phi D_\phi\bm u$, se obtiene exactamente la expresión solicitada.

Para la segunda parte partimos de la ecuación de Euler
\begin{equation}
\frac{\partial\bm u}{\partial t}+(\bm u\cdot\nabla)\bm u
=\bm g-\frac1\rho\nabla p.
\end{equation}
La contracción de $\bm u=u_r\hat r+u_\theta\hat\theta+u_\phi\hat\phi$ con el primer índice de $\nabla\bm u$ da las tres componentes convectivas. Por tanto,
\begin{align}
\hat r:\quad
&\frac{\partial u_r}{\partial t}
+u_r\frac{\partial u_r}{\partial r}
+\frac{u_\theta}{r}\frac{\partial u_r}{\partial\theta}
+\frac{u_\phi}{r\sin\theta}\frac{\partial u_r}{\partial\phi}
-\frac{u_\theta^2+u_\phi^2}{r}
=g_r-\frac1\rho\frac{\partial p}{\partial r},\\
\hat\theta:\quad
&\frac{\partial u_\theta}{\partial t}
+u_r\frac{\partial u_\theta}{\partial r}
+\frac{u_\theta}{r}\frac{\partial u_\theta}{\partial\theta}
+\frac{u_\phi}{r\sin\theta}\frac{\partial u_\theta}{\partial\phi}
+\frac{u_ru_\theta}{r}-\frac{u_\phi^2\cot\theta}{r}
=g_\theta-\frac{1}{\rho r}\frac{\partial p}{\partial\theta},\\
\hat\phi:\quad
&\frac{\partial u_\phi}{\partial t}
+u_r\frac{\partial u_\phi}{\partial r}
+\frac{u_\theta}{r}\frac{\partial u_\phi}{\partial\theta}
+\frac{u_\phi}{r\sin\theta}\frac{\partial u_\phi}{\partial\phi}
+\frac{u_ru_\phi}{r}+\frac{u_\theta u_\phi\cot\theta}{r}
=g_\phi-\frac{1}{\rho r\sin\theta}\frac{\partial p}{\partial\phi}.
\end{align}

\textbf{Pregunta 2 (3 pts).} Considere la atmósfera terrestre como un gas ideal en equilibrio hidrostático e isotérmico, de temperatura constante $T_0$, y lo suficientemente tenue como para despreciar su autogravitación. Suponga que la Tierra es una esfera sólida de radio $R_T$ y masa $M_T$, de modo que el gas está sometido únicamente al campo gravitacional terrestre.

\textbf{(a)} Si la presión y la densidad a nivel de la superficie terrestre son $p_0$ y $\rho_0$, respectivamente, determine los perfiles $p(r')$ y $\rho(r')$ en función de la altura $r'$ sobre la superficie terrestre, con $r=R_T+r'$.

\textbf{(b)} El número de Knudsen, $\mathrm{Kn}=\lambda/L$, permite estimar la validez de la aproximación de continuo. Use $\lambda\sim k_BT_0/(\pi d^2p)$ y $L\sim\rho/|\nabla\rho|$ para determinar hasta qué altura sigue siendo válida la descripción de medio continuo.

\textit{Solución.}

\textbf{(a)} En equilibrio hidrostático,
\begin{equation}
\nabla p=\rho\bm g,
\qquad
\bm g=-\frac{GM_T}{r^2}\hat r.
\end{equation}
Como el gas es ideal e isotérmico,
\begin{equation}
p=\frac{RT_0}{M}\rho,
\end{equation}
donde $M$ es la masa molar media. Sustituyendo en la ecuación hidrostática,
\begin{equation}
\frac{dp}{dr}=-\frac{A}{r^2}p,
\qquad
A:=\frac{MGM_T}{RT_0}.
\end{equation}
Integrando desde $r=R_T$, donde $p=p_0$, hasta $r=R_T+r'$, resulta
\begin{equation}
\ln\frac{p(r')}{p_0}=A\left(\frac{1}{R_T+r'}-\frac1{R_T}\right).
\end{equation}
Por tanto,
\begin{equation}
\boxed{p(r')=p_0\exp\left[A\left(\frac{1}{R_T+r'}-\frac1{R_T}\right)\right]}.\label{eq:p-profile}
\end{equation}
Como $p/\rho=RT_0/M$ es constante,
\begin{equation}
\boxed{\rho(r')=\rho_0\exp\left[A\left(\frac{1}{R_T+r'}-\frac1{R_T}\right)\right]}.\label{eq:rho-profile}
\end{equation}

\textbf{(b)} Derivando explícitamente el perfil de densidad \eqref{eq:rho-profile},
\begin{equation}
\frac{d\rho}{dr}=-\frac{A}{r^2}\rho,
\end{equation}
y por ello
\begin{equation}
L\sim\frac{\rho}{|\nabla\rho|}=\frac{r^2}{A}.
\end{equation}
Sustituyendo el perfil de presión \eqref{eq:p-profile} en la expresión dada para el camino libre medio,
\begin{equation}
\lambda(r')\sim\lambda_0\exp\left[A\left(\frac1{R_T}-\frac{1}{R_T+r'}\right)\right],
\qquad
\lambda_0:=\frac{k_BT_0}{\pi d^2p_0}.
\end{equation}
De este modo,
\begin{equation}
\boxed{\mathrm{Kn}(r')\sim
\frac{A\lambda_0}{(R_T+r')^2}
\exp\left[A\left(\frac1{R_T}-\frac{1}{R_T+r'}\right)\right]}.
\end{equation}
La condición de continuo es $\mathrm{Kn}\ll1$. Para localizar el borde de ese régimen introducimos una altura crítica $r'_c$ mediante $\mathrm{Kn}(r'_c)\sim1$. Esta igualdad de orden de magnitud no reemplaza la condición $\mathrm{Kn}\ll1$: sólo identifica el primer punto en que la separación entre escalas microscópica y macroscópica deja de ser grande.

En la superficie,
\begin{equation}
\mathrm{Kn}(0)\sim\frac{A\lambda_0}{R_T^2}.
\end{equation}
Si esta cantidad es mucho menor que uno, la descripción continua es válida inicialmente. Definimos entonces $r'_c$ como la \emph{primera} altura, al ascender desde la superficie, en la que $\mathrm{Kn}$ alcanza un valor de orden uno. Con esta definición, el régimen continuo conectado con la superficie corresponde a
\begin{equation}
\boxed{0\le r'<r'_c}.
\end{equation}
Si además se quiere comprobar que $\mathrm{Kn}$ crece en el intervalo considerado, basta observar que
\begin{equation}
\frac{d}{dr}\ln\mathrm{Kn}=\frac{A-2r}{r^2};
\end{equation}
por tanto, es monótona creciente mientras $A>2r$. Para parámetros atmosféricos terrestres usuales se cumple $A\gg 2R_T$; en forma puramente simbólica basta exigir $A>2(R_T+r'_c)$ en el intervalo estudiado. Bajo esa condición, $\mathrm{Kn}$ aumenta al ascender y queda justificado que el régimen continuo conectado con la superficie sea $r'<r'_c$.

La altura exacta queda determinada implícitamente por
\begin{equation}
\frac{A\lambda_0}{(R_T+r'_c)^2}
\exp\left[A\left(\frac1{R_T}-\frac{1}{R_T+r'_c}\right)\right]\sim1.
\end{equation}

Si además $r'\ll R_T$, entonces
\begin{equation}
\frac1{R_T+r'}=\frac1{R_T}-\frac{r'}{R_T^2}+O\!\left(\frac{r'^2}{R_T^3}\right).
\end{equation}
Definiendo la altura de escala
\begin{equation}
H:=\frac{R_T^2}{A}=\frac{RT_0R_T^2}{MGM_T},
\end{equation}
se obtiene
\begin{equation}
L\sim H,
\qquad
\lambda(r')\sim\lambda_0e^{r'/H},
\qquad
\mathrm{Kn}(r')\sim\frac{\lambda_0}{H}e^{r'/H}.
\end{equation}
Así,
\begin{equation}
\boxed{r'_c\sim H\ln\left(\frac{H}{\lambda_0}\right)}.
\end{equation}
La expansión usada es autoconsistente sólo si el resultado verifica
$\,r'_c/R_T=(H/R_T)\ln(H/\lambda_0)\ll1\,$. Nótese que $H=R_T^2/A$ tiene dimensiones de longitud: no escala simplemente como $R_T^2$, porque $A$ también tiene dimensiones de longitud. Esta condición debe comprobarse a posteriori.

\end{document}
